\begin{table}[htbp]
    \caption{Negative-Weight Diagnostic for the Baseline TWFE Design}
    \label{tab:negative-weights}
    \centering
    \small
    \renewcommand{\arraystretch}{1.08}
    \resizebox{0.92\textwidth}{!}{%
    \begin{tabular}{lcccc}
        \doubletoprule
        Outcome & Negative components & Total components & Share negative & Negative weight mass \\
        \midrule
        Log registered voters & 3 & 5,460 & 0.055\% & 0.000004 \\
        Low-education share & 3 & 5,460 & 0.055\% & 0.000004 \\
        High-education share & 3 & 5,460 & 0.055\% & 0.000004 \\
        \doublebottomrule
    \end{tabular}}
    \vspace{1.0em}
    \begin{minipage}{0.92\textwidth}
        {\footnotesize \textbf{Note:} The table reports the de Chaisemartin-D'Haultfoeuille negative-weight diagnostic using \texttt{TwoWayFEWeights} for the absorbing-treatment version of the baseline BVR design. The reported share negative is the fraction of ATT components receiving negative weight under the implied two-way fixed-effects aggregation; the final column reports the absolute mass of those negative weights. \par}
    \end{minipage}
\end{table}
